% ============================================================================
% TP5 -- Paper, version 1 (entregable OE5)
% Formato IEEE de congreso (IEEEtran), dos columnas, 8 carillas, en ingles.
% Revision doble ciego: el autor va anonimo hasta la version final.
% Plantilla armada a partir del paper de referencia de la catedra (solo el
% formato y la estructura, sin su texto).
% Compilar con pdfLaTeX (Overleaf: Menu > Compiler > pdfLaTeX).
% ============================================================================
\documentclass[conference]{IEEEtran}

\usepackage[utf8]{inputenc}
\usepackage[T1]{fontenc}
\usepackage{cite}
\usepackage{amsmath,amssymb}
\usepackage{graphicx}
\usepackage{booktabs}
\usepackage{xcolor}
\usepackage[hidelinks]{hyperref}

\begin{document}

\title{Software Defined Radio-Based Reference Source\\
for EMC Radiated Emissions (RE) Verification and\\
Calibration}

% Version anonima para la revision. En la version final van los autores:
% \author{\IEEEauthorblockN{Martín Ledesma}
% \IEEEauthorblockA{Facultad de Ingeniería, Universidad de Buenos Aires\\
% Buenos Aires, Argentina\\
% maledesma@fi.uba.ar}}
\author{\IEEEauthorblockN{Anonymous Submission for Double-Blind Review}
\IEEEauthorblockA{Author information withheld for peer review}}

\maketitle

\begin{abstract}
Abstract text (one paragraph, about 150 to 250 words): problem, proposed system, calibration methodology, experimental validation and main results.
\end{abstract}

\begin{IEEEkeywords}
electromagnetic compatibility (EMC), radiated emissions, software defined radio (SDR), comb generator, reference source, calibration.
\end{IEEEkeywords}

% ============================================================================
\section{Introduction}

Context and motivation. Problem statement. Contribution of this work.

The main contributions of this work are:
\begin{enumerate}
  \item First contribution.
  \item Second contribution.
  \item Third contribution.
\end{enumerate}

The remainder of this paper is organized as follows. Section~\ref{sec:state} reviews the state of the art. Sections~\ref{sec:proposed} and~\ref{sec:implementation} describe the proposed system and its implementation. Section~\ref{sec:results} presents the experimental results and Section~\ref{sec:emissions} the emissions measurements. Sections~\ref{sec:conclusions} and~\ref{sec:future} present the conclusions and future work.

% ============================================================================
\section{State of the Art}\label{sec:state}

Review of previous work, with references in IEEE style~\cite{b1}, \cite{b2}.

\begin{table}[htbp]
  \caption{Comparison with representative implementations}
  \label{tab:comparison}
  \centering
  \begin{tabular}{@{}lcccc@{}}
    \toprule
    Work & Tech. & Port. & Prog. & EMC \\
    \midrule
    Author (year) & -- & -- & -- & -- \\
    Author (year) & -- & -- & -- & -- \\
    This work     & SDR & Yes & Yes & Yes \\
    \bottomrule
  \end{tabular}
\end{table}

Table~\ref{tab:comparison} summarizes representative implementations.

% ============================================================================
\section{Proposed System}\label{sec:proposed}

Architecture and design objectives.

\begin{figure}[htbp]
  \centering
  \fbox{\parbox[c][3cm][c]{0.9\columnwidth}{\centering Figure placeholder\\ (replace with \texttt{\textbackslash includegraphics})}}
  \caption{Functional block diagram of the proposed system.}
  \label{fig:block}
\end{figure}

Fig.~\ref{fig:block} presents the block diagram of the system.

% ============================================================================
\section{System Implementation}\label{sec:implementation}

Hardware, software and calibration.

\begin{equation}
  E_{dB\mu V/m} = V_{receiver} + AF + L_{cable} - G_{amp}
  \label{eq:field}
\end{equation}

% ============================================================================
\section{Experimental Results}\label{sec:results}

\subsection{First Analysis}
Text.

\subsection{Second Analysis}
Text.

% ============================================================================
\section{Emissions Measurements}\label{sec:emissions}

\subsection{Conducted Emissions}
Text.

\subsection{Radiated Emissions}
Text.

% ============================================================================
\section{Conclusions}\label{sec:conclusions}

Text.

% ============================================================================
\section{Future Work}\label{sec:future}

Text.

% ============================================================================
\begin{thebibliography}{00}
\bibitem{b1} A. Author, B. Author, and C. Author, ``Title of the paper,'' \textit{Journal Name}, vol. 1, no. 1, pp. 1--10, 2026.
\bibitem{b2} A. Author, \textit{Title of the Book}. City: Publisher, 2026.
\end{thebibliography}

\end{document}
